\documentclass[a0paper,portrait]{baposter}
\usepackage[french]{babel}


\usepackage{wrapfig}
\usepackage{lmodern}
\usepackage{lipsum,graphicx}
\usepackage{fontspec}
\usepackage{advdate}
\usepackage{svg}
\usepackage{xcolor}
\usepackage{tabularx}
\usepackage{enumitem}


% Set main font (Abraham)
% \setmainfont[
%     Path = font/Abraham/,
%     Extension = .otf,
%     % SmallCapsFont = *-smallcaps,
%     BoldFont = *-bold, % Outline
%     ItalicFont = *-italic, % Italic
%     BoldItalicFont = *-bolditalic % Outline Italic
% ]{abraham}
% \setmainfont{TeX Gyre Termes}

%

\graphicspath{{figures/}} % Directory in which figures are stored

\newcommand{\compresslist}{%
\setlength{\itemsep}{0pt}%
\setlength{\parskip}{1pt}%
\setlength{\parsep}{0pt}%
}

\newenvironment{boenumerate}
  {\begin{enumerate}\renewcommand\labelenumi{\textbf\theenumi.}}
  {\end{enumerate}}


\begin{document}

\definecolor{Mycolor1}{HTML}{3FFF00}
\definecolor{Mycolor2}{HTML}{0FF000} %0FF000

\begin{poster}
{
grid=false,
headerborder=open, % Adds a border around the header of content boxes
colspacing=1em, % Column spacing
bgColorOne=white, % Background color for the gradient on the left side of the poster
bgColorTwo=white, % Background color for the gradient on the right side of the poster
borderColor=Mycolor1, % Border color
headerColorOne=Mycolor2, % Background color for the header in the content boxes (left side)
headerColorTwo=Mycolor2, % Background color for the header in the content boxes (right side)
headerFontColor=white, % Text color for the header text in the content boxes
boxColorOne=white, % Background color of the content boxes
textborder=rounded, %rectangle, % Format of the border around content boxes, can be: none, bars, coils, triangles, rectangle, rounded, roundedsmall, roundedright or faded
eyecatcher=false, % Set to false for ignoring the left logo in the title and move the title left
headerheight=0.08\textheight, % Height of the header
headershape=rounded, % Specify the rounded corner in the content box headers, can be: rectangle, small-rounded, roundedright, roundedleft or rounded
headershade=plain,
headerfont=\Large\textsf, % Large, bold and sans serif font in the headers of content boxes
%textfont={\setlength{\parindent}{1.5em}}, % Uncomment for paragraph indentation
linewidth=1pt % Width of the border lines around content boxes
}
{}
%
%----------------------------------------------------------------------------------------
%	TITLE AND AUTHOR NAME
%----------------------------------------------------------------------------------------
%
{\textsf{{Tableau de bord - Éco-marathon Shell 2023}}} % Titre en Sans Serif
{\sf\vspace{0.1em}\\
Électrique  -  {\AdvanceDate[0]\today} 
\vspace{0.1em}\\
\small{ Période couverte du {\AdvanceDate[-7]\today} au {\AdvanceDate[0]\today}
}
}
{\includegraphics[width=.1\linewidth]{img/udes.pdf}} % TU Dublin logo

\headerbox{Objectif de projet}{name=objectif,column=0,row=0, span=3}{
Créer un véhicule efficace à propulsion électrique pour gagner la compétition de l'Éco-marathon Shell dans la division ``véhicule urbain'' en 2023.

\vspace{0.2cm} %vertical spacing
}


\headerbox{Objectifs de session}{name=objectif_session,column=1,below=objectif,span=2}{

%\vspace{0.15cm}
\centering
{    \begin{tabularx}{\linewidth}{
        >{\hsize=1.5\hsize}X
        >{\hsize=1.0\hsize}X
        >{\centering\arraybackslash\hsize=0.5\hsize}X
      }
        
        \textbf{Objectif} & \textbf{Système} & \textbf{Avancement} \\
         Modélisation numérique & Moteur/Transmission & 100\% \\
         CAD détaillé & Moteur/Transmission & 99\% \\
         Assemblage V0 & Moteur/Transmission & 32\% \\
         PCB V2 onduleur & Onduleur & 75\% \\
         PCB de contrôle fonctionnel & Instrumentation & 30\% \\
         BMS V2 commandé & Bloc batterie & 60\% \\
         Bloc batterie V0 assemblé et testé & Bloc batterie & 85\% \\
         PCBs de lumière fonctionnels et assemblés & Instrumentation & 60\% \\
         Essuie & Instrumentation & 40\% \\
      \end{tabularx}
    
        
     % Template des Objectifs rencontés:
     %
     % Objectif : l'objectif de la session
     % Système  : le nom ou numéro système que l'objectif touche ex: Simulateur ou SIM1
     % &  & = ESPACE POUR LE FORMATAGE NE PAS ENLEVER
     % Avancement : Pourcentage estimé d'avancement de l'objectif.
     %
     %  
     %  \textbf{Objectif} & \textbf{Système} &  & \textbf{Avancement}\\
     %   Objectif & Système &  & Avancement \\
     %   Objectif & Système &  & Avancement \\
     %   Objectif & Système &  & Avancement \\
     %   Objectif & Système &  & Avancement \\
     %   Objectif & Système &  & Avancement \\
    
    
    
    
    
}

% \vspace{0.4cm} % spacing vertical

}

\headerbox{Ordre du jour}{name=ordre_jour,column=0,below=objectif,span=1}{
\vspace{0.1cm}
\begin{itemize}
    \item \textbf{Rencontre gestion :}
    \begin{enumerate}
        \item Revue conception et commande PCB V1 onduleur
        \item Revue de conception et commande PCB V1 BMS
        \item Changement de plan DCDC 12V
        %%\item Quatrième truc à parler gestion
        %%\item Cinquième truc à parler gestion
    \end{enumerate}
\end{itemize}
\vspace{0.25cm} % spacing vertical
}

\headerbox{Membres}{name=membres,column=2,below=ordre_jour,span=1}{
\vspace{0.3cm}


% inserts an image inside the box, 5 rows high | r = on the right along side text 
\begin{wrapfigure}[5]{r}{0.25\textwidth}
\includegraphics[width=0.9\linewidth]{img/membres/François-Paquette-3.jpg} 
\end{wrapfigure}
\subsubsection*{}
\vspace{-2mm}
\textbf{François Paquette}

Batterie / BMS

Moteur

\begin{wrapfigure}[5]{r}{0.25\textwidth}
\includegraphics[width=.9\linewidth]{img/membres/Gabriel-Cabana-2.jpg} 
\end{wrapfigure}
\subsubsection*{}
\vspace{-2mm}
\textbf{Gabriel Cabana}

Moteur

Instrumentation


\begin{wrapfigure}[5]{r}{0.25\textwidth}
\includegraphics[width=.9\linewidth]{img/membres/Jérôme-Gelé-2.jpg} 
\end{wrapfigure}
\subsubsection*{}
\vspace{2mm}
\textbf{Jérôme Gelé}

Contrôle moteur \& Onduleur

Batterie / BMS

\begin{wrapfigure}[5]{r}{0.25\textwidth}
\includegraphics[width=.9\linewidth]{img/membres/Joël-Grégoire-Lagueux-1.jpg} 
\end{wrapfigure}
\subsubsection*{}
\vspace{2mm}
\textbf{Joël Grégoire-Lagueux}

Instrumentation

\begin{wrapfigure}[5]{r}{0.25\textwidth}
\includegraphics[width=.9\linewidth]{img/membres/Loïc-Poirier-2.jpg} 
\end{wrapfigure}
\subsubsection*{}
\vspace{-2mm}
\textbf{Loïc Poirier}

Batterie / BMS

Simulateur

\begin{wrapfigure}[5]{r}{0.25\textwidth}
\includegraphics[width=.9\linewidth]{img/membres/Marc-Antoine-Dubreuil-2.jpg} 
\end{wrapfigure}
\subsubsection*{}
\vspace{-2mm}
\textbf{Marc-Antoine Dubreuil}

Contrôle moteur \& Onduleur

Simulateur

\begin{wrapfigure}[5]{r}{0.25\textwidth}
\includegraphics[width=.9\linewidth]{img/membres/Thomas-Chagnon-3.jpg} 
\end{wrapfigure}
\subsubsection*{}
\vspace{2mm}
\textbf{Thomas Chagnon}

Moteur / Transmission

\begin{wrapfigure}[5]{r}{0.25\textwidth}
\includegraphics[width=.9\linewidth]{img/membres/Vincent-Bonneau-3.jpg} 
\end{wrapfigure}
\subsubsection*{}
\vspace{2mm}
\textbf{Vincent Bonneau}

Contrôle moteur \& Onduleur


\begin{wrapfigure}[5]{r}{0.25\textwidth}
\includegraphics[width=.9\linewidth]{img/membres/Xavier-Morin-2.jpg} 
\end{wrapfigure}
\subsubsection*{}
\vspace{2mm}
\textbf{Xavier Morin}

Moteur / Transmission

Transferts thermiques


\vspace{1cm}

}


\headerbox{Courbe en S}{name=courbe_s,column=0,span=2,below=ordre_jour}{
{
    {
        \hspace{0.75cm}
        \centering
    	\includegraphics[width=0.9\linewidth]{img/Courbe_S.pdf}
    }
}

}

\headerbox{Heures travaillées}{name=heures,column=0,span=2,below=courbe_s}{
{
    
    {
        \hspace{0.75cm}
        \centering
    	\includegraphics[width=0.9\linewidth]{img/Heures_travaillees.pdf}
    }
    % \vspace{1.3cm}
}

}

\end{poster}

\begin{poster}
{
grid=false,
headerborder=open, % Adds a border around the header of content boxes
colspacing=1em, % Column spacing
bgColorOne=white, % Background color for the gradient on the left side of the poster
bgColorTwo=white, % Background color for the gradient on the right side of the poster
borderColor=Mycolor1, % Border color
headerColorOne=Mycolor2, % Background color for the header in the content boxes (left side)
headerColorTwo=Mycolor2, % Background color for the header in the content boxes (right side)
headerFontColor=white, % Text color for the header text in the content boxes
boxColorOne=white, % Background color of the content boxes
textborder=rounded, %rectangle, % Format of the border around content boxes, can be: none, bars, coils, triangles, rectangle, rounded, roundedsmall, roundedright or faded
eyecatcher=false, % Set to false for ignoring the left logo in the title and move the title left
headerheight=0\textheight, % Height of the header
headershape=rounded, % Specify the rounded corner in the content box headers, can be: rectangle, small-rounded, roundedright, roundedleft or rounded
headershade=plain,
headerfont=\Large\textsf, % Large, bold and sans serif font in the headers of content boxes
%textfont={\setlength{\parindent}{1.5em}}, % Uncomment for paragraph indentation
linewidth=1pt % Width of the border lines around content boxes
}
{}
{\textsf{{}}} % Titre vide pcq sinon ça plante
\hfill \break % Ligne vide pcq sinon ça plante

\headerbox{Objectifs bimensuels}{name=objectifs,column=0, row=0, span=3}{{\begin{tabularx}{\linewidth}{
    |>{\hsize=2.25\hsize}X|% 10% of 4\hsize 
    >{\hsize=0.75\hsize}X|% 30% of 4\hsize
    >{\hsize=0.75\hsize}X|% 30% of 4\hsize 
    >{\centering\arraybackslash\hsize=0.25\hsize}X|% 30% of 4\hsize 
       % sum=0.2\hsize for 4 columns
  }
    \hline
    \textbf{Tâche} & \textbf{Système} & \textbf{Responsable} & \textbf{Heures}\\\hline
     Modélisation & Thermique & Xavier Morin & 4\\\hline
     Modélisation et calculs & Transmission & Xavier Morin & 9\\\hline
     Steering Committee & Gestion & Xavier Morin & 2\\\hline
     Choix du circuit intégré du système de gestion des batteries & Bloc batterie / BMS & Loïc Poirier & 4\\\hline
     Début de la conception du convertisseur CC/CC & Bloc batterie / BMS & Vincent Bonneau & 4\\\hline
     Lire la documentation sur le contrôle moteur & Contrôle moteur & Jérôme Gelé & 2\\\hline
     Choix du circuit intégré du système de gestion des batteries & BMS & Jérôme Gelé & 2\\\hline
     Conception préliminaire du circuit & Bloc batterie / BMS & Loïc Poirier & 4\\\hline
     Modélisation du contrôle moteur & Contrôle moteur & Vincent Bonneau & 4\\\hline
    
     Schématique de l'onduleur et librairies Altium & Contrôle moteur & Marc-Antoine Dubreuil & 6\\\hline
     Modélisation préliminaire du moteur en accélération & Moteur & Gabriel Cabana & 12\\\hline
  \end{tabularx}

% TEMPLATE des lignes du tableau de taches accomplis cette semaine
%
% Tâche: résumé de la tâche
% Système : Numéro de système ou nom complet ex : Simulateur ou SIM1
% Responsable : Initiales ou nom complet ex : Gabriel Cabana ou G.C.
% Heures : Heures passées par le responsable à faire la tâche la semaine dernière (jeudi à jeudi)
%
%\textbf{Tâche} & \textbf{Système} & \textbf{Responsable} & \textbf{Heures}\\\hline
%    Tâche & Système & Responsable & Heures\\\hline
%    Tâche & Système & Responsable & Heures\\\hline
%    Tâche & Système & Responsable & Heures\\\hline
%    Tâche & Système & Responsable & Heures\\\hline
%    Tâche & Système & Responsable & Heures\\\hline
%    Tâche & Système & Responsable & Heures\\\hline
%    Tâche & Système & Responsable & Heures\\\hline
%    Tâche & Système & Responsable & Heures\\\hline
%    Tâche & Système & Responsable & Heures\\\hline
%    Tâche & Système & Responsable & Heures\\\hline
%    Tâche & Système & Responsable & Heures\\\hline
%    Tâche & Système & Responsable & Heures\\\hline
%  }}


\headerbox{Problèmes}{name=problemes,column=0,span=3,below=objectifs}{
\begin{tabularx}{\linewidth}{
    |>{\hsize=1.7\hsize}X|% 10% of 3\hsize 
    >{\hsize=0.5\hsize}X|% 30% of 3\hsize
    >{\hsize=0.8\hsize}X|% 30% of 3\hsize 
       % sum=0.2\hsize for 4 columns
  }
    \hline
    \textbf{Problème} & \textbf{Système} & \textbf{Responsable} \\\hline
    Quincaillerie manquante J.M. Bergeron & Transmission & Xavier Morin et Thomas Chagnon\\\hline
    Assemblage soudé arbres des roues & Transmission & Xavier Morin et Thomas Chagnon\\\hline
    Assemblage logement à roulements et longeron véhicule & Transmission & Xavier Morin et Thomas Chagnon\\\hline
    Mauvaise organisation de la documentation PCB de contrôle & Instrumentation & Vincent Bonneau\\\hline
\end{tabularx}
    
    
 % Template des problèmes rencontés:
 %
 % Problème : le problème rencontré cette semaine
 % Système  : le nom ou numéro système que le problème touche ex: Simulateur ou SIM1
 % Responsable : Nom ou Initiales du ou des personnes touchées par ce problèmes ex: G.C. C.E.G.
 %
 %  \hline
 %  \textbf{Problème} & \textbf{Système} & \textbf{Responsable}\\\hline
 %  Problème & Système & Responsable \\\hline
 %  Problème & Système & Responsable \\\hline
 %  Problème & Système & Responsable \\\hline
 %  Problème & Système & Responsable \\\hline
 %  Problème & Système & Responsable \\\hline
 %  Problème & Système & Responsable \\\hline
 %  Problème & Système & Responsable \\\hline
 %
}

\headerbox{Risques}{name=risques,column=0,span=3,below=problemes}{
\begin{tabularx}{\linewidth}{
    |>{\let\newline\\\hsize=0.40\hsize}X|% 10% of 4\hsize 
    >{\hsize=0.25\hsize}X|% 30% of 4\hsize
    >{\hsize=0.25\hsize}X|% 30% of 4\hsize 
    >{\centering\arraybackslash\hsize=0.1\hsize}X|% 30% of 4\hsize 
       % sum=0.2\hsize for 4 columns
  }
    \hline
    \textbf{Risque} & \textbf{Prévention} & \textbf{Gestion} & \textbf{Priorité}\\\hline
  % Mauvaises connexion des phases du moteur & revues des pairs avant exécution & Rebobinage du moteur ou utilisation du moteur de remplacement & HAUTE\\\hline
  
    Moteur V0 non fonctionnel & Mettre les bouchées double afin d'avoir un prototype & Se limiter à un moteur acheté ou faire usiner à l'extérieur & HAUTE\\\hline
    
    Moteur V0 ou acheté non caractérisé pour l'échéancier du 6 octobre & Intégration partielle au véhicule & Communication active avec les équipes interdépendantes & MOYENNE\\\hline
    
    %Moteur qui ne respecte pas les spécifications acceptables & Faire simulations, calculs, revues de conception avec experts & Achat d'un moteur plan B avec date limite & HAUTE\\\hline
    

    Mauvaise conception système onduleur & Revue de conception avec experts, validation V1, validation PCB contrôle, attribuer MVP & acquisition onduleur (att. Budget/Temps) & HAUTE\\\hline
    
    Mauvaise conception du PCB de lumière & Système simple et revue de conception inter-équipe & Modifications directes sur le PCB ou achat de composants aditionnels au besoin & HAUTE\\\hline
    
    Commande des connecteurs tardive & Communication directe avec commanditaire & Commande faite par l'équipe au lieu du commanditaire & MOYENNE\\\hline
    

  \end{tabularx}
  
  
% Template pour le tableau des risques de la semaine : 

% Risque : le risque de la semaine
% Mitigation: comment mitiger le risque pour réduire ses conséquences/occurances/etc.
% Conséquence : La ou les conséquences si le risque survient
% Priorité : Priorité du risque sur une échelle de 1 à 5.  5 étant + prioritaire. La priorité est basé sur les conséquences la probablité d'occurence, etc.  
  
%  \textbf{Risque} & \textbf{Mitigation} & \textbf{Conséquence} & \textbf{Priorité}\\\hline
%    Risque & Mitigation & Conséquence & Priorité\\\hline
%    Risque & Mitigation & Conséquence & Priorité\\\hline
%    Risque & Mitigation & Conséquence & Priorité\\\hline
%    Risque & Mitigation & Conséquence & Priorité\\\hline
%    Risque & Mitigation & Conséquence & Priorité\\\hline
%    Risque & Mitigation & Conséquence & Priorité\\\hline
%    Risque & Mitigation & Conséquence & Priorité\\\hline
%    Risque & Mitigation & Conséquence & Priorité\\\hline
%    Risque & Mitigation & Conséquence & Priorité\\\hline
%    Risque & Mitigation & Conséquence & Priorité\\\hline
%    Risque & Mitigation & Conséquence & Priorité\\\hline
%    Risque & Mitigation & Conséquence & Priorité\\\hline
}

\end{poster}

%%%%%%%%%%%%%%%%%%%%%%%%%%%%%%%%%%%%%%%%%%%%%%%%%%

\begin{poster}
{
grid=false,
headerborder=open, % Adds a border around the header of content boxes
colspacing=1em, % Column spacing
bgColorOne=white, % Background color for the gradient on the left side of the poster
bgColorTwo=white, % Background color for the gradient on the right side of the poster
borderColor=Mycolor1, % Border color
headerColorOne=Mycolor2, % Background color for the header in the content boxes (left side)
headerColorTwo=Mycolor2, % Background color for the header in the content boxes (right side)
headerFontColor=white, % Text color for the header text in the content boxes
boxColorOne=white, % Background color of the content boxes
textborder=rounded, %rectangle, % Format of the border around content boxes, can be: none, bars, coils, triangles, rectangle, rounded, roundedsmall, roundedright or faded
eyecatcher=false, % Set to false for ignoring the left logo in the title and move the title left
headerheight=0\textheight, % Height of the header
headershape=rounded, % Specify the rounded corner in the content box headers, can be: rectangle, small-rounded, roundedright, roundedleft or rounded
headershade=plain,
headerfont=\Large\textsf, % Large, bold and sans serif font in the headers of content boxes
%textfont={\setlength{\parindent}{1.5em}}, % Uncomment for paragraph indentation
linewidth=1pt % Width of the border lines around content boxes
}
{}
{\textsf{{}}} % Titre vide pcq sinon ça plante
\hfill \break % Ligne vide pcq sinon ça plante

\headerbox{Cashflow}{name=cashflow,column=0,span=3}{
% \vspace{.6cm} % spacing vertical
\hspace{1cm}
\begin{center}
    \includegraphics[width=\linewidth]{img/Cashflow_text.pdf}
\end{center}
}

\headerbox{Citation de la semaine}{name=citation,column=0,row=0, span=3,below=cashflow}{
 \textit{"Le travail fatigue même les ânes."} -- Mateo Aleman
%
\vspace{0.2cm} %vertical spacing
}

\end{poster}

\end{document}